\section{Exact certificate for the three-star bound}\label{app:tape}

The certificate establishes~\eqref{eq:norm-test} on the
nondominant part of the sorted nonnegative cone in six coordinates.
It covers each of the fifteen perfect matchings. Homogeneity permits a
unit-sum slice for ray computations. The numerical search that proposed
the subdivision and active sets supplies no inequality or coverage step.
The input to exact verification is the saved finite plan and the explicit
geometry below.

\subsection{Polyhedral partition}

Sort $x_1\ge\cdots\ge x_6\ge0$ and retain all index matchings,
including when some values tie or vanish. For a matching
$\mathcal M$ let $P_\mathcal M x$ consist of $x_i+kx_j$ for each
pair $i<j$. Then $B=\norm{P_\mathcal M x}$ and
\[
 A(x)=\frac1{32}\sum_{s\in\{1\}\times\{-1,1\}^5}\abs{s^\top x}.
\]
Hausmann and Rodriguez \cite[Table (1--1), Sections 4--5]{hausmann}
classify the same length arrangement. Each wall equates the lengths
in a subset with those in its complement. Writing their difference
as a sign form, reversing the sorting convention and taking closures
identifies the walls here. Their chamber counts in six and eight
coordinates are 21 and 2,470, respectively. The exact verification
below reconstructs the six-coordinate geometry and uses it for a
matching-dependent inequality.

In gap variables $d_i=x_i-x_{i+1}$ with $x_7=0$, the coefficients of
$s^\top x$ are the successive prefix sums of $s$. Exactly twelve
forms have both positive and negative gap coefficients and need
both orientations. Every other form has a constant weak sign.
For each of the $2^{12}$ weak orientations, the intersection with
the nonnegative gap cone is pointed. Every extreme ray has five
linearly independent active homogeneous normals, with defining
equalities also available on lower-dimensional chambers.
Enumerate all five-normal choices from the six gap facets and
twelve walls and take their signed cofactors. Retain nonzero
nonnegative null vectors, convert gaps to coefficients, and normalize
each to primitive integer coordinates. For each orientation retain
all rays satisfying its weak inequalities. This procedure lists
every extreme ray and hence generates every closed cone.
Dropping duplicate ray sets and proper subsets preserves the union.
The exact result is eighteen rays and twenty-one maximal chambers.
Each maximal cone has six independent generators, checked by its
determinant. The one cone with strict dominance in its interior is
handled by Lemma~\ref{lem:dominant}. Its shared faces remain covered.
Zero weights and ties belong to closed chambers, and the origin
gives zero inner moments.

\subsection{A certificate on one closed cone}

Throughout one chamber $A=\ell^\top x$ with $\ell$ a signed
average of sign vectors. Then $|\ell^\top z|\le A(z)\le\norm z$
in the entire ambient space. Put $N=I-\ell\ell^\top$. It is positive
definite on each nondominant maximal chamber. Indeed,
$\norm\ell=1$ would give $A(\ell)=1$, so the absolute Rademacher
sum at $\ell$ would be constant by Cauchy--Schwarz. The degree-two
sign coefficients of its square imply $\ell_i\ell_j=0$ for $i\ne j$.
Positivity and $A(x)\ge x_1$ on a full-dimensional sorted chamber
then force $\ell=e_1$. Averaging first over the largest sign gives
$A(x)=x_1$ exactly in the dominant region.

For generators $v_i$ normalized to unit coordinate sum put
$G_{ij}=v_i^\top N v_j$, and choose
$\beta_i\ge \norm{P_\mathcal M v_i}-A(v_i)$. For a certificate
vector $h$ verify
\begin{equation}
 Gh\ge\beta,\qquad h^\top Gh\le H,\quad H=4-8\sqrt2/3 .
 \label{eq:check}
\end{equation}
With $x=\sum_i t_i v_i$ for $t_i\ge0$ and $y=\sum_i h_i v_i$,
the triangle inequality,~\eqref{eq:check} and Cauchy--Schwarz
in $N$ give
\[
 B(x)-A(x)\le\sum_i t_i\beta_i
 \le x^\top Ny\le \sqrt{x^\top Nx}\sqrt{h^\top Gh}.
\]
This is the required continuous inequality. The entries of $h$
may have either sign. At a stored active set $J$ the checker sets
$h_{J^c}=0$, solves $G_{JJ}h_J=\beta_J$ and checks every row of
$Gh\ge\beta$. The norm equals $h_J^\top\beta_J$.

A ray norm is exact for zero, one nonzero entry or two equal nonzero
entries of $P_\mathcal M v_i$. Otherwise the least upper bound on
the grid of denominator $10^{12}$ is used, obtained by comparing
squares. The ray norm is at most one by unit-sum normalization.
All inequalities use $\mathbb Q(\sqrt2)$. For opposite signs of
the two rational parts of $a+b\sqrt2$, testing its sign reduces
to comparing $a^2$ and $2b^2$ and selecting the dominant sign.
The linear solves have rational Gram matrices and use positive
pivots. All field comparisons are exact.

\subsection{Subdivision plan and verification interface}

Sort the primitive integer rays lexicographically. A chamber is
identified by its increasing tuple of ray indices in that list.
Visit chambers lexicographically and omit the dominant chamber.
Within each chamber visit matchings in the recursive order that
pairs the least unused index with each remaining index in increasing
order and recurses on the others. This yields $20\cdot15=300$
root cases. A base64 block in the preserved S07 document decodes
to the preorder traversal of those binary trees.
Bytes $1$ through $63$ are nonempty active-set masks on the six
current generators. Bytes $64$ through $78$ select one of the
fifteen pairs of generator slots in lexicographic order.

For a split on slots $i<j$, replace $v_j$ by $(v_i+v_j)/2$ in
the first child and $v_i$ by the same midpoint in the second.
Both children remain simplicial. If $t_i\ge t_j$, use $2t_j$
on the midpoint and $t_i-t_j$ on $v_i$ in the first child.
If $t_i\le t_j$ use the reversed representation in the second.
They cover the entire parent and share the equality face. The
verifier visits both branches and consumes the plan exactly.
It checks 1,854 nodes, including 777 splits and 1,077 terminal
certificates, at a maximum depth of six. The node count also satisfies $1854=2(1077)-300$.

On a strict norm leaf,~\eqref{eq:check} gives a strict bound for
all nonzero $x$ by positive definiteness of $N$. Exactly 34 leaves
saturate the norm condition. For every such leaf the verifier
checks that $y$ is a positive multiple of $(x_0,x_0,1,1,0,0)$.
Equality in Cauchy--Schwarz forces $x$ onto this ray. This supplies
a necessary equality direction on every terminal leaf. Sharpness
itself follows from the explicit attaining weights in the main
proof.

The exact plan, embedded in the S07 copy, and the two Python
verifiers accompany the manuscript under \path{certificate/}.
See \path{README.md} for the two replay commands. S07 enumerates
rays by cofactors and bounds the norm by field comparisons. S08
independently reconstructs the partition by clipping the sorted
simplex against the walls, including all opposite-side intersections
and filtering by active-normal rank. It uses full sign sums,
an alternative field representation, integer square-root brackets
and rational $LDL^\top$ solves. The two implementations verify the
same proposed subdivision and active-set data. S08's optional
eight-coordinate geometry run counts 2,470 maximal cells, 206
rays and a maximum of 37 rays in a cell. Its checks concern these
partition statistics, with shared faces included by the geometry.
A certificate of all four-star inequalities would require separate
checks on the remaining domain.
